\section{Deck Strategy and Interpretability}

Alakazam's Powerful Hand provides an unusually interpretable local objective:
two damage counters are placed for each card in hand. For a visible target with
remaining HP $H$, the basic hand threshold is
\begin{equation}
  h_{\mathrm{KO}}=\left\lceil H/20\right\rceil,
\end{equation}
modified only by visible, applicable effects. Public deck guides provide useful
context~\cite{tcgplayerAlakazamGuide,ptcgoAlakazamGuide}, but the project guide
is training curriculum metadata, not a serving action source. It cannot read
hidden facts, provide policy cross-entropy, or override the learned option
score.

The strategy is resource-constrained card-count control, not unconstrained
draw. A stable board maintains a current attacker route, a replacement Abra or
evolved line, a Dunsparce draw/recycle route when it solves a defined problem,
and a flexible bench slot. Deterministic search normally precedes broad draw
when the target is already known. Rare Candy is preserved when evolving through
Kadabra keeps multiple routes live. Boss's Orders is spent for a prize-changing
or tempo-winning target; Enhanced Hammer is used when a visible Special Energy
actually prevents damage or changes the prize race. Near deck-out, mandatory
turn draw and every unavoidable attachment, evolution, or ability draw are
counted before committing the line.

The 14 August matchup supplement illustrates how semantic features connect to
human reasoning without introducing hidden inference. A visible 330-HP Mega
Lopunny ex requires 17 hand cards for a full-health Powerful Hand knockout and
gives up three Prizes. A visibly attached Mist Energy prevents the attack effect,
so Enhanced Hammer or a gust-around line is valuable. The controller may use
visible HP, card class, attachment, and legal options. It may not infer a hidden
Mist Energy or healing card merely because the opponent revealed Buneary.

Interpretability follows the same causal boundary. For each current option the
trace can expose the base logit, structured rule residual, masks, bounded route
reliabilities, and final score. Leave-one-route-out logit changes measure
decision sensitivity, not counterfactual win effect. A route should be
reweighted only after paired fixed-seed intervention. Likewise, checklist
records explain what evidence was available, but their exact-zero gates at this
snapshot mean they did not change the deployed action.
